\documentclass[crop,tikz,11pt]{standalone}
\usepackage{geometricfigures}
\usepackage{amsmath,amssymb}
\usepackage{pgfplots}
\pgfplotsset{compat=1.18}
\usetikzlibrary{arrows.meta,decorations.markings,patterns,intersections,calc}
\usepgfplotslibrary{fillbetween}

\definecolor{lib}{HTML}{3B75AF}
\definecolor{rot}{HTML}{C1701E}
\definecolor{sep}{HTML}{B03030}
\definecolor{net}{HTML}{4E9A4E}

% One set of curves on the cylinder, drawn through three charts, all symplectic and all on one
% scale, so areas can be compared by eye across the panels.
%   (a) the angular strip (theta, p), theta in radians;
%   (b) the round chart with kappa = 1:
%         (theta, p) -> rho (-cos theta, -sin theta),  rho^2/2 = a - p;
%   (c) (b) followed by the two shears (x, y) -> (x + s y^2, y) -> (x, y + t x^2), each with
%       Jacobian determinant 1, so (c) is symplectic as well.
% The curves: librating H = 0.5 and H = 0.95 (bounce points at theta = arccos H), the separatrix,
% the zero section p = 0 (the midline, dotted), and the rotating level H = 1.6 on both branches,
% p < 0 solid, p > 0 dashed. The grey disk is inside the innermost curve drawn.
\pgfmathsetmacro{\a}{3.2}
\pgfmathsetmacro{\sx}{0.16}
\pgfmathsetmacro{\sy}{0.13}

% the p of each family as a function of the angle t (degrees)
\pgfmathdeclarefunction{pl}{2}{\pgfmathparse{sqrt(max(0,2*((#2)-cos(#1))))}}
\pgfmathdeclarefunction{ps}{1}{\pgfmathparse{2*abs(sin((#1)/2))}}
% chart (a)
\pgfmathdeclarefunction{ax}{2}{\pgfmathparse{rad(#1)}}
\pgfmathdeclarefunction{ay}{2}{\pgfmathparse{#2}}
% chart (b)
\pgfmathdeclarefunction{bx}{2}{\pgfmathparse{-sqrt(2*(\a-(#2)))*cos(#1)}}
\pgfmathdeclarefunction{by}{2}{\pgfmathparse{-sqrt(2*(\a-(#2)))*sin(#1)}}
% chart (c): u = bx + sx by^2, then (u, by + sy u^2)
\pgfmathdeclarefunction{cx}{2}{\pgfmathparse{bx(#1,#2)+\sx*by(#1,#2)*by(#1,#2)}}
\pgfmathdeclarefunction{cy}{2}{\pgfmathparse{by(#1,#2)+\sy*cx(#1,#2)*cx(#1,#2)}}

% \hole{x-chart}{y-chart}: the grey region inside the innermost curve drawn, the p > 0 rotating
% orbit at H = 1.6; only the planar charts have one
\newcommand{\hole}[2]{%
  \addplot[draw=none,fill=fg!8!bg,domain=0:360] ({#1(x,pl(x,1.6))},{#2(x,pl(x,1.6))}) -- cycle;}
% \orbits{x-chart}{y-chart}: every curve of the figure through one chart
\newcommand{\orbits}[2]{%
  % librating orbits: the p > 0 leg, then the p < 0 leg back
  \foreach \H/\ta/\tb in {0.5/60/300, 0.95/18.19/341.81}{
    \addplot[lib,thick,domain=\ta:\tb,samples=200]
      ({#1(x,pl(x,\H))},{#2(x,pl(x,\H))});
    \addplot[lib,thick,domain=\ta:\tb,samples=200]
      ({#1(x,-pl(x,\H))},{#2(x,-pl(x,\H))});
  }
  % the midline: the zero section, through both equilibria
  \addplot[fg!70!bg,densely dotted,thick,domain=0:360] ({#1(x,0)},{#2(x,0)});
  % separatrix branches
  \addplot[sep,very thick,domain=0:360] ({#1(x,ps(x))},{#2(x,ps(x))});
  \addplot[sep,very thick,domain=0:360] ({#1(x,-ps(x))},{#2(x,-ps(x))});
  % rotating orbits at H = 1.6
  \addplot[rot,thick,domain=0:360] ({#1(x,-pl(x,1.6))},{#2(x,-pl(x,1.6))});
  \addplot[rot,thick,dashed,domain=0:360] ({#1(x,pl(x,1.6))},{#2(x,pl(x,1.6))});
  % bounce points of the two librating orbits, on the midline
  \foreach \t in {60,300,18.19,341.81}{
    \addplot[only marks,mark=*,mark size=1.4pt,lib] coordinates {({#1(\t,0)},{#2(\t,0)})};
  }
  % the stable equilibrium (theta = pi) and the unstable one P (theta = 0)
  \addplot[only marks,mark=*,mark size=1.5pt,fg] coordinates {({#1(180,0)},{#2(180,0)})};
  \addplot[only marks,mark=*,mark size=1.9pt,sep] coordinates {({#1(0,0)},{#2(0,0)})};
}

\begin{document}
\begin{tikzpicture}
\pgfplotsset{every axis/.style={x=0.62cm,y=0.62cm,hide axis,clip=false},
  every axis plot/.append style={smooth,samples=240,variable=x}}

% ---------------------------------------------------------------- (a)
\begin{axis}[name=A,xmin=-0.1,xmax=6.4,ymin=-3.7,ymax=3.7]
  \orbits{ax}{ay}
  % the seam: theta = 0 and theta = 2 pi are one line of the cylinder
  \draw[fg!35!bg] (axis cs:0,-3.1) -- (axis cs:0,3.1);
  \draw[fg!35!bg] (axis cs:{2*pi},-3.1) -- (axis cs:{2*pi},3.1);
  \draw[fg!55!bg,thick] (axis cs:-0.12,2.75) -- (axis cs:0,2.9) -- (axis cs:0.12,2.75);
  \draw[fg!55!bg,thick] (axis cs:{2*pi-0.12},2.75) -- (axis cs:{2*pi},2.9) -- (axis cs:{2*pi+0.12},2.75);
  \node[sep,font=\scriptsize,left] at (axis cs:0,0) {$P$};
  \node[sep,font=\scriptsize,right] at (axis cs:{2*pi},0) {$P$};
  \node[font=\scriptsize,below] at (axis cs:0,-3.1) {$0$};
  \node[font=\scriptsize,below] at (axis cs:pi,-3.1) {$\theta=\pi$};
  \node[font=\scriptsize,below] at (axis cs:{2*pi},-3.1) {$2\pi$};
  \node[font=\scriptsize,sep,fill=bg,inner sep=1pt] at (axis cs:4.95,1.25) {$\ell_+$};
  \node[font=\scriptsize,sep,fill=bg,inner sep=1pt] at (axis cs:4.95,-1.25) {$\ell_-$};
  \node[font=\scriptsize,fg!70!bg,fill=bg,inner sep=1pt] at (axis cs:2.3,0.3) {$\nu$};
  \node[font=\footnotesize,anchor=north] at (axis cs:pi,-3.95) {(a) the cylinder, cut open};
\end{axis}

% ---------------------------------------------------------------- (b)
\begin{axis}[name=B,at={($(A.east)+(0.7cm,0)$)},anchor=west,
  xmin=-3.6,xmax=3.6,ymin=-3.7,ymax=3.7]
  \hole{bx}{by}
  \orbits{bx}{by}
  \node[sep,font=\scriptsize,left] at (axis cs:{-sqrt(2*\a)},0) {$P$};
  \node[font=\scriptsize,fg!60!bg] at (axis cs:0,0) {$K$};
  \node[font=\footnotesize,anchor=north] at (axis cs:0,-3.95) {(b) round chart, $\rho^2/2=a-p_\theta$};
\end{axis}

% ---------------------------------------------------------------- (c)
\begin{axis}[name=C,at={($(B.east)+(0.5cm,0)$)},anchor=west,
  xmin=-2.6,xmax=4.4,ymin=-3.7,ymax=3.7]
  \hole{cx}{cy}
  \orbits{cx}{cy}
  \node[sep,font=\scriptsize,left] at (axis cs:{cx(0,0)},{cy(0,0)}) {$P$};
  \node[font=\scriptsize,fg!60!bg] at (axis cs:0.25,0.1) {$K$};
  \node[font=\footnotesize,anchor=north] at (axis cs:0.9,-3.95) {(c) (b), sheared};
\end{axis}
\end{tikzpicture}
\end{document}
